\section{Trainability and Stability: Effective Rank and the Sign Failure}\label{sec:trainability}

\paragraph{Trainability is a second axis.}
Proposition~1 guarantees expressiveness, not trainability. At complete rank every basis
represents the target exactly, yet on M4-Hourly the DCT and the Karhunen--Lo\"eve basis fail to
train on every seed under a plain linear head (Table~\ref{tab:effective_rank}). The obstacle is
not the basis matrix---every basis here is orthonormal, so its condition number is one---but the
coefficient vector the head must produce. M4-Hourly windows carry $0.98$ of their energy in a
single direction, and each of the two failing bases aligns one row with it, so the head must
learn one coefficient near $225{,}000$ beside forty-seven near $240$ from one linear layer at
one learning rate. The effective rank---the participation ratio $1/\sum_k e_k^2$ of the target's
coefficient energy shares, computable from the basis and the data before any training---measures
exactly this concentration: $1.05$ and $1.54$ for the two bases that fail, $17.86$ and $17.12$
for the two that never do. Across bases it tracks the measured convergence rate at
$\rho = +0.949$, where captured energy, flat to the fourth decimal, explains nothing.

\begin{table}[t]
  \centering
  \caption{M4-Hourly at complete forecast rank ($\df = \Ho = 48$), plain linear synthesis head,
  ten seeds. Effective rank is the participation ratio of the target's coefficient spectrum. The
  two concentrated bases fail on every seed; the two diffuse bases converge on every seed.
  Effective rank tracks the convergence rate across bases at $\rho = +0.949$; captured energy is
  flat to the fourth decimal across these bases and explains nothing.}
  \label{tab:effective_rank}
  \begin{tabular}{lcc}
    \toprule
    Basis & Effective rank & Converged seeds (of 10) \\
    \midrule
    Karhunen--Lo\"eve   & $1.05$  & 0 \\
    DCT                 & $1.54$  & 0 \\
    Wavelet (band-ordered) & $17.86$ & 10 \\
    Fourier-style       & $17.12$ & 10 \\
    \bottomrule
  \end{tabular}
\end{table}

\paragraph{The trend-branch repair.}
A paired polynomial trend branch absorbs the dominant component before the frozen basis sees it,
and repairs the failure without adding expressive power: the DCT goes from $28.748$ to $9.111$
\smape{} and Karhunen--Lo\"eve from $27.607$ to $9.291$, each from zero of ten seeds to ten of
ten ($p = 0.00018$; a control without the trend branch remains at zero of ten). This is why the
completeness null of Section~\ref{sec:completeness} is a trend-branch-present null: every arm of
the swap grid carried the repair, so the null speaks only to configurations that have it.

\paragraph{The divergence mechanism.}
The instability that remains is not a general loss of fit. The model occasionally predicts a
negative value where the target must be positive. \smape{} is
$200\,|\hat{y}-y|/(|\hat{y}|+|y|)$, and opposite signs give $|\hat{y}-y| = |\hat{y}|+|y|$, so
such a point scores exactly $200$ whatever its magnitude. Saving every forecast from forty runs
shows that this is the whole effect (Table~\ref{tab:sign_failure}): diverged runs put between
$7.1$ and $7.8$ percent of their points below zero, the fraction scoring $200$ equals the
fraction below zero to every digit, and deleting exactly those points recovers the converged
score to within $0.11$ \smape.

\begin{table}[t]
  \centering
  \caption{Every forecast saved from forty runs. Divergence on M4 is a sign failure, not a fit
  failure: the points at the \smape{} ceiling of $200$ are exactly the points below zero, and
  removing them recovers the converged score.}
  \label{tab:sign_failure}
  \begin{tabular}{@{}ll@{}}
    \toprule
    Measurement & Diverged runs \\
    \midrule
    Forecast points below zero & $7.1$--$7.8$\% \\
    Points at the \smape{} ceiling of $200$ & equals the below-zero fraction, to every digit \\
    Score with the below-zero points deleted & within $0.11$ \smape{} of the converged score \\
    \bottomrule
  \end{tabular}
\end{table}

\paragraph{Remedies.}
Bottlenecks, frozen low-rank bases, backcast penalties, and output rectifiers all remove the
mode, because all constrain what a block can inject; none adds expressive power. A rectifier
removes the mode by construction on non-negative targets, though on sub-band decompositions it
destroys the cancellation between signed corrections and is catastrophic---the rule is
architecture-scoped, not universal. The $\dagger$ mechanism (the trunk activation applied to the
forecast head; Table~\ref{tab:legend}) eliminates the stuck-mode failure entirely where
tested---six-stack Milk from $63\%$ to $0\%$ of runs, ten-stack Milk from $64\%$ to $0\%$, and
M4-Yearly from $4.5\%$ to $0\%$---but it costs accuracy on the converged subset on every M4
period except Hourly. We treat it as a targeted intervention, not a global default (full table
in Appendix~\ref{sec:ablations}).

\paragraph{Two quantities, computable before training.}
Captured energy (Section~\ref{sec:truncation}) says how accurate a truncated basis will be;
effective rank says whether it will train at all. Both are functions of the basis and the data
alone, and neither is a property of the wavelet family.
